% Perfil de lectura digital durante la redacción.
% Evita la alternancia de márgenes y la corrección de encuadernación de 10 mm
% que la clase oficial activa para impresión a doble cara.
\KOMAoptions{open=any,BCOR=0mm}
\recalctypearea

% Capítulo, sección, subsección y subsubsección en el índice general.
\setcounter{tocdepth}{3}

% Margen de maniobra para identificadores técnicos en párrafos estrechos.
% Evita desbordamientos mínimos sin reducir la tipografía ni alterar tablas.
\setlength{\emergencystretch}{3em}

% Los identificadores técnicos de las tablas (dominios, rutas, hashes, etc.)
% necesitan poder partirse aunque no contengan espacios. El filtro de Quarto
% breakable-table-code.lua aplica esta macro únicamente al código en línea
% situado dentro de tablas.
\usepackage{seqsplit}
\newcommand{\BreakableTableCode}[1]{\texttt{\seqsplit{#1}}}

% La clase oficial antepone "Figura" y "Tabla" al número de cada entrada,
% pero su ancho dinámico solo mide el número (por ejemplo, 4.1). Reservamos
% explícitamente espacio para la etiqueta completa y evitamos solapamientos.
\DeclareTOCStyleEntry[
  numwidth=7em,
  entrynumberformat=\useheaderfont\figurename~\autodot,
  linefill=\hfill,
  entryformat=\uselabelfont,
  pagenumberformat=\useheaderfont
]{tocline}{figure}
\DeclareTOCStyleEntry[
  numwidth=7em,
  entrynumberformat=\useheaderfont\tablename~\autodot,
  linefill=\hfill,
  entryformat=\uselabelfont,
  pagenumberformat=\useheaderfont
]{tocline}{table}

% Compatibilidad entre la línea decorativa de etsibook y las cajas de título
% que genera Quarto. Se conserva el aspecto oficial usando el ancho efectivo
% de la caja, en lugar del ancho global de página.
\makeatletter
\renewcommand*{\chapterlinesformat}[3]{%
  \if\relax\detokenize{#2}\relax
    {\raggedleft#3\par}%
    {\noindent\color{gray!75}\rule{\linewidth}{3.5pt}\par}%
  \else
    \@hangfrom{#2}{#3}%
    \par\nobreak
    {\noindent\color{gray!75}\rule{\linewidth}{3.5pt}\par}%
  \fi
}
\makeatother
